\resumeSubHeadingListStart
    \resumeSubheading
    {iOS Developer, iOS Team}{Jan 2019 -- Oct 2021}
    {}{강남, 서울}

    \resumeProjectHeading
    {\textbf{iOS 어플리케이션 개발 및 배포}}{iOS Developer}
    \resumeItemListStart
        \resumeItem{\textbf{핵심 문제 및 목표: } 미술 관련 컨텐츠를 네이티브 앱에서 }
        \resumeItem{\textbf{나의 역할: } 네이티브 iOS 어플리케이션의 API Client를 포함한 Backend Layer를 주도하여 구축. 작품, 작가, 전시 상세 View 및 ViewController 구현, Unity Bridge 를 이용하여 Google Cardboard VR 기능 구현, Google Cartdboard VR 기능을 위한 Asset들을 다운로드 할 수 있도록 Download 기능 구현, ARKit을 이용하여 Artwork Hanger 기능 구현}
        \resumeItem{\textbf{문제 해결 과정:}}
            \resumeItemListStart
                \resumeItem{\textbf{기술적 고민 및 선택 (Why):} VR 에셋 다운로드, Unity 연동 등 복잡한 기능 추가를 앞두고, 기존 MVC 구조의 복잡성 증가 및 유지보수성 저하를 예상. 상태(State)를 보다 명확하고 예측 가능하게 관리할 수 있는 반응형 아키텍처 도입을 목표로 함.}
                \resumeItem{\textbf{트레이드오프 분석 (Trade-off):} \textbf{직접 MVVM 패턴을 구현}하는 방안을 먼저 고려함. 하지만 이는 상태와 사용자 입력을 처리하기 위한 프로토콜 정의, 데이터 바인딩 로직 구현 등 상당한 초기 개발 리소스가 필요하다고 판단. 대안으로, 반응형 프로그래밍 프레임워크인 \textbf{ReactorKit}을 검토함. ReactorKit은 \textbf{사용자 입력(Action) -> 비즈니스 로직 처리(Mutate) -> 상태 변경 적용(Reduce) }으로 이어지는 3단계의 명확한 데이터 흐름을 강제하여, 직접 구현 시 발생할 수 있는 복잡성을 원천적으로 차단하고 코드의 흐름을 매우 쉽게 이해할 수 있도록 만들어주는 점이 인상적이었음. 초기 학습 곡선을 감수하더라도, ReactorKit을 도입하는 것이 장기적인 유지보수성과 개발 효율성 측면에서 더 유리하다고 결론 내림.}
                \resumeItem{\textbf{구체적인 구현 내용 (How):} ReactorKit을 기반으로 앱의 아키텍처를 재설계. 모든 사용자 입력은 'Action'으로, 상태 변화는 'State'로 정의하여 데이터 흐름을 통일하고, AWS S3 클라이언트를 이용한 다운로드 서비스를 Reactor에 통합하여 구현.}
            \resumeItemListEnd
         \resumeItem{\textbf{결과 및 임팩트:}}
            \resumeItemListStart
                \resumeItem{\textbf{개발 효율성 향상:} 프레임워크에 익숙해진 후, \textbf{Action-Mutate-Reduce 구조에 맞춰 로직을 정의하는 과정이 매우 쉬워지고 명확해짐.} 이로 인해 신규 기능 추가 시 개발 속도가 향상되고 코드의 일관성이 유지됨.}
                \resumeItem{\textbf{안정성 및 유지보수성 확보:} 명확한 구조 덕분에 상태 관리 관련 버그 발생률이 감소했으며, 새로운 팀원이 코드 구조를 파악하는 시간이 단축됨.}
            \resumeItemListEnd
    \resumeItemListEnd
\resumeSubHeadingListEnd